\documentclass{claudeeditorial}

\renewcommand{\DocKicker}{Technical Showcase}
\renewcommand{\DocTitle}{Documentación técnica editorial}
\renewcommand{\DocSubtitle}{Código, SQL, tablas, notas laterales y diagramas con una estética cálida y reproducible}
\renewcommand{\DocAuthor}{Equipo de documentación}
\renewcommand{\DocRole}{Plantilla pública}
\renewcommand{\DocDate}{5 de mayo de 2026}
\renewcommand{\DocRef}{TECH-001}
\renewcommand{\DocFooter}{claudeeditorial technical showcase}

\hypersetup{
  pdftitle={Claude Editorial Technical Showcase},
  pdfauthor={Equipo de documentación}
}

\ClaudeRunningHeader
\setstretch{1.2}

\begin{document}
\BodyCopy

\ClaudeCover

\ClaudeBanner[pulseDecision]{Código, datos y arquitectura en un mismo sistema visual}

\begin{tcbraster}[raster columns=3,raster equal height=rows,raster column skip=8pt]
  \KPICard{Código}{12+}{lenguajes frecuentes}
  \KPICard{Diagramas}{UML/ER}{TikZ nativo}
  \KPICard{Pandoc}{Ready}{tokens Highlighting}
\end{tcbraster}

\section{Bloques semánticos}

\ClaudeMarginInfo[Nota lateral]{Las notas marginales son compactas; funcionan mejor con una o dos frases.}

\begin{ClaudeInfo}[Contexto]
Este documento muestra piezas técnicas neutras. La clase evita depender de paquetes UML/ER externos y ofrece estilos TikZ propios para mantener una licencia permisiva.
\end{ClaudeInfo}

\begin{ClaudeTip}[Criterio editorial]
Usa cajas semánticas para explicar intención: información, advertencia, validación, pregunta o riesgo. No las uses como decoración.
\end{ClaudeTip}

\begin{tcbraster}[raster columns=2,raster equal height=rows,raster column skip=10pt]
\begin{ClaudeWarning}[Riesgo de maquetación]
Los listados largos deben partir líneas y reservar espacio suficiente a la numeración.
\end{ClaudeWarning}

\begin{ClaudeCheck}[Validación]
Las tablas técnicas usan color alterno, cabeceras cálidas y reglas suaves.
\end{ClaudeCheck}
\end{tcbraster}

\section{Código}

\begin{ClaudeCode}[title={Consulta SQL con alias y ordenación}]{SQL}
SELECT
  customer_id,
  COUNT(*) AS order_count,
  SUM(total_amount) AS total_revenue
FROM analytics.orders
WHERE status = 'paid'
GROUP BY customer_id
HAVING SUM(total_amount) > 1000
ORDER BY total_revenue DESC;
\end{ClaudeCode}

\begin{ClaudeCode}[title={Servicio TypeScript}]{TypeScript}
type Decision = {
  id: string;
  confidence: number;
  accepted: boolean;
};

export function selectDecision(items: Decision[]): Decision | undefined {
  return items
    .filter((item) => item.accepted)
    .sort((a, b) => b.confidence - a.confidence)[0];
}
\end{ClaudeCode}

\begin{ClaudeCode}[title={Clases Java}]{Java}
interface Humano {
  void identificate();
}

class Persona implements Humano {
  private String nombre;

  @Override
  public void identificate() {
    System.out.println("Persona: " + nombre);
  }
}
\end{ClaudeCode}

\begin{ClaudeCode}[title={Modelo C\#}]{CSharp}
public abstract class Volumen
{
    public string Nombre { get; set; }
    public int Numero { get; set; }

    public abstract void Mostrar();
}
\end{ClaudeCode}

\begin{ClaudeConsole}
$ xelatex template.tex
Output written on template.pdf
\end{ClaudeConsole}

\section{Tablas}

\begin{ClaudeTable}{@{}lLr@{}}
\ClaudeTableHeader
Elemento & Uso recomendado & Prioridad \\
\toprule
Portada & Presentar título, subtítulo y metadatos con aire editorial. & 1 \\
Info box & Explicar contexto sin interrumpir el flujo principal. & 2 \\
Diagrama & Mostrar estructura cuando el texto se vuelve demasiado lineal. & 1 \\
Código & Documentar implementación, consulta o salida reproducible. & 1 \\
\bottomrule
\end{ClaudeTable}

\section{Diagramas}

\begin{ClaudeDiagram}
\centering
\begin{tikzpicture}[claude diagram canvas,node distance=22mm]
  \ClaudeUMLInterface{Repository}{}{+ find(id) : Entity\\+ save(entity) : void}
  \ClaudeUMLClass[right=of Repository]{Service}
    {- repository : Repository}
    {+ execute(command) : Result\\- validate(command) : bool}
  \ClaudeUMLClass[right=of Service]{Controller}
    {- service : Service}
    {+ handle(request) : Response}
  \ClaudeUMLRelation[claude uml dependency]{Service}{Repository}[uses]
  \ClaudeUMLRelation[claude uml relation]{Controller}{Service}[calls]
\end{tikzpicture}
\end{ClaudeDiagram}

\begin{ClaudeDiagram}
\centering
\begin{tikzpicture}[claude diagram canvas,node distance=18mm]
  \ClaudeEREntity{customer}{Customer}
  \ClaudeERAttribute[above left=of customer]{customer_id}{\ClaudeKey{id}}
  \ClaudeERAttribute[above right=of customer]{customer_name}{name}
  \ClaudeERRelationship[right=of customer]{places}{places}
  \ClaudeEREntity[right=of places]{order}{Order}
  \ClaudeERAttribute[above=of order]{order_id}{\ClaudeKey{id}}
  \ClaudeERAttribute[below=of order]{order_total}{total}
  \ClaudeEROneToMany{customer}{places}[1..N]
  \ClaudeERManyToOne{places}{order}[N..1]
  \ClaudeERLink{customer_id}{customer}
  \ClaudeERLink{customer_name}{customer}
  \ClaudeERLink{order_id}{order}
  \ClaudeERLink{order_total}{order}
\end{tikzpicture}
\end{ClaudeDiagram}

\ClaudeSignature{\DocAuthor}{\DocRole}

\end{document}
